%% kapa: exam=1 round=21 year=2010 period=1 form=A subject=economics subject_ko=경제원론 seq=10 number=50
%% source: Q-Net 감정평가사 기출문제 (cst00309) artlSeq=5024181 file=extracted/1008/5024181/2010년 제21회 감정평가사 1차시험 기출문제/1교시.hwp
%% conversion: kordoc 4.24.51 -> tools/gen1.py (auto); review=proofread-sonnet; flags=none
\begin{question}{50}
자유변동환율제도(free floating exchange rate system)에 관한 설명으로 옳지 않은 것은?
\begin{choices}
  \item 고정환율제도에 비해서 상대적으로 통화정책의 자주성을 확보할 수 있다.
  \item 환율의 신속한 시장수급 조절기능은 대외 균형을 유지하는데 도움이 된다.
  \item 환율변동에 따른 교역당사자의 환위험 부담이 있다.
  \item 각국의 정책당국들이 경쟁적으로 평가절상정책을 실시한다.
  \item 각국의 이자율 수준이 환율 결정에 영향을 미친다.
\end{choices}
\end{question}
\answer{4}
